\scanpage{20}
en
\[
0\unsure{-|\to}S(A)-STS(A)-S\unsure{Z|\zeta}(X)\longrightarrow0
\]

elle est exacte en $STS(A)$ et $SZ(X)$ puisque $S$ est exact à droite, et exacte à gauche et \unsure{splittée} car on a \unsure{explicitmt|explicitement|explicite} un inverse à \deleted{droite} gauche de $S(A)\to S(TS)(A)=(ST)S(A)$.

[ Soit $M\in\mathcal C$ \unsure{telle|tel} que pour \deleted{\unsure{tt|toute}} suite \deleted{de la} $S$-exacte de la forme $\Sigma(A)$, la suite transformée $\mathrm{Hom}(\Sigma(A),M))$ soit exacte :
\[
0-\mathrm{Hom}(Z(A),M)-\mathrm{Hom}(TS(A),M)-\mathrm{Hom}(A,M)\to0
\]

Alors $M$ est $S$-injectif. En effet, on fera $M=A$, et on trouve que $A-TS(A)$ identifie $A$ à un \uline{facteur direct} de $TS(A)=T(X)$ $(X=S(A))$. Or $T(X)$ est $S$-injectif en vertu de (i), donc \deleted{tout} facteur direct l'est, donc $A$ l'est. On voit en \unsure{$=$} temps que $M$ $S$-injectif $\Longleftrightarrow$ $M\to TS(M)$ est un isomorphisme sur un facteur direct (cf les faiblement injectifs de Cartan Eilenberg).

\uline{Remarque} Supposons \deleted{que le foncteur $S$ soit exact} \inserted{le foncteur $S$ exact \deleted{à droite}} \deleted{et} que $S(u)=0$ implique $u=0$, \inserted{i.e. $\varphi$ est \uline{injectif}.} Alors une suite $S$-exacte \inserted{\unsure{(1)|($\Sigma$)}} est exacte. En effet, \deleted{pour \unsure{tout} $A\in\mathcal C$} \deleted{$S$\unsure{pour} \illegible{1}} si $\Sigma:P\xrightarrow{u}Q\overset{v}{-}R$ et $S(v)S(u)=0$, \deleted{\unsure{alors|donc}} \inserted{i.e.} $S(vu)=0$, alors

% archive sheet number: 19
